% !TeX root = ../main.tex
\section{Теоретические сведения}

\subsection{Основные параметры полосовых фильтров}
В радиочастотных и СВЧ-трактах фильтры служат для селекции полезного сигнала и подавления внеполосных помех. Поведение линейного четырёхполюсника описывается матрицей рассеяния ($S$-параметрами), где коэффициент передачи задаётся комплексной величиной:
\begin{equation}
    S_{21}(\omega) = |S_{21}(\omega)| e^{j\varphi(\omega)},
\end{equation}
где $|S_{21}(\omega)|$ определяет амплитудно-частотную характеристику (АЧХ), а $\varphi(\omega) = \arg S_{21}(\omega)$ — фазо-частотную характеристику (ФЧХ).

К основным параметрам фильтра относятся:
\begin{itemize}
    \item \textbf{Центральная частота $f_0$} — средняя частота полосы пропускания;
    \item \textbf{Полоса пропускания по уровню $-3\text{ дБ}$ ($\Delta f_{-3\text{ дБ}}$)}:
    \begin{equation}
        \Delta f_{-3\text{ дБ}} = f_{\text{в}} - f_{\text{н}},
    \end{equation}
    где $f_{\text{н}}$ и $f_{\text{в}}$ — нижняя и верхняя граничные частоты среза по уровню половинной мощности;
    \item \textbf{Вносимое затухание ($IL$)} на центральной частоте:
    \begin{equation}
        IL = -20 \lg |S_{21}(f_0)| \quad [\text{дБ}];
    \end{equation}
    \item \textbf{Неравномерность в полосе пропускания ($\delta L$)} — разность между максимальным и минимальным затуханием в рабочей полосе;
    \item \textbf{Групповое время запаздывания (ГВЗ, $\tau_{\text{гр}}$)} — скорость изменения фазового сдвига по частоте:
    \begin{equation}
        \tau_{\text{гр}}(\omega) = -\frac{d\varphi(\omega)}{d\omega}.
    \end{equation}
\end{itemize}

\begin{figure}[htbp]
    \centering
    \includegraphics[width=1\textwidth]{ris8.png}
    \caption{Типовая амплитудно-частотная характеристика полосового фильтра}
    \label{fig:ris8}
\end{figure}

\subsection{Коэффициент передачи и схема его измерения}
Для измерения АЧХ исследуемого фильтра используется метод калибровки сквозного тракта с последующим включением четырёхполюсника в разрыв кабелей. Генератор сигналов формирует зондирующий синусоидальный сигнал заданной частоты и мощности, а анализатор спектра фиксирует прошедшую мощность.

\begin{figure}[htbp]
    \centering
    \includegraphics[width=1\textwidth]{ris10.png}
    \caption{Структурная схема измерения коэффициента передачи фильтра}
    \label{fig:ris10}
\end{figure}

Коэффициент передачи фильтра на каждой частоте определяется разностью показаний анализатора спектра при наличии фильтра $P_{\text{изм}}(f)$ и без него (режим калибровки тракта) $P_{\text{кал}}(f)$:
\begin{equation}
    |S_{21}(f)|_{\text{дБ}} = P_{\text{изм}}(f) - P_{\text{кал}}(f).
\end{equation}

\subsection{Коэффициент отражения, КСВН и методика их измерения}
При подключении к тракту с волновым сопротивлением $Z_0 = 50\text{ Ом}$ нагрузки $Z_L$ модуль коэффициента отражения по напряжению $\Gamma$ и коэффициент стоячей волны по напряжению (КСВН) связаны выражениями[cite: 1, 2]:
\begin{equation}
    \Gamma = \frac{Z_L - Z_0}{Z_L + Z_0}, \qquad \text{КСВН} = \frac{1 + |\Gamma|}{1 - |\Gamma|}.
\end{equation}

В логарифмическом масштабе модуль коэффициента отражения (Return Loss) определяется как:
\begin{equation}
    |\Gamma|_{\text{дБ}} = 20 \lg |\Gamma|.
\end{equation}

Для разделения падающей и отражённой волн применяется направленный ответвитель, на вход которого (\textit{IN}) подаётся сигнал генератора, к выходу (\textit{OUT}) подключается исследуемый объект, а сигнал с плеча связи (\textit{CPL OUT}) поступает на вход анализатора спектра[cite: 2].

\begin{figure}[htbp]
    \centering
    % Вставьте файл с рисунком 13 из методических указаний
    \includegraphics[width=1\textwidth]{ris13.png}
    \caption{Структурная схема измерения модуля коэффициента отражения помощью направленного ответвителя}
    \label{fig:ris13}
\end{figure}

Калибровка схемы производится подключением холостого хода (разомкнутая цепь, $|\Gamma| = 1 \implies 0\text{ дБ}$) к порту \textit{OUT} с регистрацией мощности $P_{\text{open}}(f)$[cite: 1, 2]. Затем к порту подключается фильтр, нагруженный на согласованную нагрузку $50\text{ Ом}$, и измеряется мощность $P_{\text{filt}}(f)$[cite: 1, 2]. Коэффициент отражения фильтра вычисляется по формуле[cite: 1]:
\begin{equation}
    |\Gamma_{\text{фильтр}}(f)|_{\text{дБ}} = P_{\text{filt}}(f) - P_{\text{open}}(f).
\end{equation}
